% fig_lec04_coord_basis.tex — Coordinate basis diagram for Lecture 4
% Data: generated by scripts/sim_lec03_charts.jl
%
% Shows f : O ⊂ M → R   and   ψ : O → U ⊂ R²,
% so that  f∘ψ⁻¹ : U → R   is a function on Euclidean space
% whose partial derivatives at ψ(p) define the coordinate basis vectors.

\begin{center}
\begin{tikzpicture}
  \pgfplotsset{
    colormap={cgbluetint}{rgb255(0cm)=(210,224,240) rgb255(1cm)=(112,156,200)},
  }

  % ── TOP: M with O_α and p ─────────────────────────────────
  \begin{axis}[
      name=M,
      width=7.2cm, height=5.3cm,
      view={-28}{35},
      xmin=-3, xmax=3, ymin=-3, ymax=3, zmin=-1.1, zmax=1.1,
      axis lines=none, xtick=\empty, ytick=\empty, ztick=\empty, clip=false,
  ]
    \addplot3[surf, opacity=0.28, shader=interp,
              colormap name=grsurface, mesh/rows=55, forget plot]
      table {data/lec03_charts_surface.dat};
    \addplot3[surf, opacity=0.55, shader=interp,
              colormap name=cgbluetint, mesh/rows=18, forget plot]
      table {data/lec03_charts_Oalpha_patch.dat};
    \addplot3[only marks, mark=*, mark size=1.6pt, banana!90!black, forget plot]
      table {data/lec03_charts_point_p_M.dat};
    \node[font=\sf\small, text=spacecadet] at (axis cs:-2.4,-2.8,1.0) {$\M$};
    \node[font=\small, text=cgblue] at (axis cs:0.2,0.8,0.75) {$O$};
    \node[font=\small, text=banana!60!black] at (axis cs:1.25,-0.1,0.32) {$p$};
  \end{axis}

  % ── BOTTOM LEFT: target R of f ────────────────────────────
  \begin{axis}[
      name=Rtarget,
      at={($(M.south)+(-3.0cm,-1.1cm)$)}, anchor=north,
      width=3.4cm, height=1.8cm,
      axis lines=middle,
      axis line style={birrenorange, line width=0.6pt,
                       -{Stealth[length=5pt,width=3pt]}},
      xmin=-1.2, xmax=1.6, ymin=-0.3, ymax=0.3,
      xtick=\empty, ytick=\empty,
      xlabel={}, ylabel={},
      clip=false,
  ]
    \addplot[only marks, mark=*, mark size=1.5pt, banana!90!black, forget plot]
      coordinates {(0.45,0)};
    \node[font=\sf\small, text=birrenorange] at (axis cs:1.45,0.18) {$\R$};
    \node[font=\scriptsize, text=banana!60!black, anchor=south]
      at (axis cs:0.45,0.04) {$f(p)$};
  \end{axis}

  % ── BOTTOM RIGHT: U ⊂ R² with ψ(p) ────────────────────────
  \begin{axis}[
      name=U,
      at={($(M.south)+(3.0cm,-1.1cm)$)}, anchor=north,
      width=4.4cm, height=4.4cm,
      axis lines=middle,
      axis line style={birrenorange, line width=0.6pt,
                       -{Stealth[length=5pt,width=3pt]}},
      xmin=-1.4, xmax=2.6, ymin=-1.1, ymax=1.8,
      xtick=\empty, ytick=\empty, clip=false,
  ]
    \addplot[fill=cgblue!15, draw=none, forget plot]
      table {data/lec03_charts_Ualpha_boundary.dat} \closedcycle;
    \addplot[only marks, mark=*, mark size=1.5pt, banana!90!black, forget plot]
      table {data/lec03_charts_point_p_alpha.dat};
    \node[font=\sf\small, text=birrenorange] at (axis cs:2.4,1.55) {$\Rn{n}$};
    \node[font=\small, text=cgblue] at (axis cs:0.5,0.3) {$U$};
    \node[font=\scriptsize, text=banana!60!black, anchor=north]
      at (axis cs:0.95,-0.14) {$\psi(p)$};
  \end{axis}

  % ── Arrows ─────────────────────────────────────────────────
  % f : M → R (down-left)
  \draw[thick, -{Stealth[length=6pt,width=3.5pt]}, spacecadet]
    ($(M.south)+(-0.3,-0.05)$) to[out=-110,in=70]
    ($(Rtarget.north east)+(-0.15,0.05)$)
    node[midway, left, font=\small, text=spacecadet, xshift=-1pt]{$f$};

  % ψ : M → R^n (down-right)
  \draw[thick, -{Stealth[length=6pt,width=3.5pt]}, cgblue]
    ($(M.south)+(0.3,-0.05)$) to[out=-70,in=110]
    ($(U.north west)+(0.15,0.05)$)
    node[midway, right, font=\small, text=cgblue, xshift=1pt]{$\psi$};

  % composite f ∘ ψ⁻¹ from U back to R
  \draw[thick, -{Stealth[length=6pt,width=3.5pt]}, spacecadet]
    ($(U.west)+(-0.05,0)$) -- ($(Rtarget.east)+(0.05,0)$)
    node[midway, below=1pt, font=\scriptsize, text=spacecadet]
      {$f\circ\psi^{-1}$};
\end{tikzpicture}
\end{center}
